\section{Android-Images}

Android nutzt GPT mit vielen Partitionen.
Wichtigste Partition: USERDATA (enthält App-Daten, Fotos, Datenbanken).

\begin{lstlisting}[language=bash]
mmls android.E01
# Falls Speicherzugriffsfehler:
xmount --in ewf android.E01 \
  --cache /tmp/phone.ovl --out raw /tmp/phone
mmls /tmp/phone/android.dd

losetup --partscan --find --show /tmp/phone/android.dd
ll /dev/loop*
\end{lstlisting}

\section{SQLite-Datenbanken}

Android speichert fast alles in SQLite: Kontakte, SMS, Anruflisten, Browser-History, GPS.

\begin{lstlisting}[language=bash]
# Alle Datenbanken finden
fls -pr /dev/loopXpN | grep '\.db$'

# Nur nicht-geloeschte Datenbanken
fls -pr /dev/loopXpN | grep '\.db$' | grep -v '*'

# Datenbank extrahieren und oeffnen
icat /dev/loopXpN <INODE> > /tmp/database.db
file /tmp/database.db
sqlitebrowser /tmp/database.db
# oder:
sqlite3 /tmp/database.db ".tables"
\end{lstlisting}

Wichtig: WAL-Dateien (\texttt{.db-wal}) können aktuellere Daten enthalten als die Hauptdatenbank.
Autopsy liest WAL automatisch mit.

\section{EXIF/GPS-Daten}

Fotos enthalten oft GPS-Koordinaten in EXIF-Metadaten:

\begin{lstlisting}[language=bash]
# Alle Bilder mit GPS-Daten auslesen
exiftool -csv -GPSLatitude -GPSLongitude \
  -GPSLatitudeRef -GPSLongitudeRef \
  -DateTimeOriginal /tmp/photos/*.jpg > /tmp/gps_data.csv

# Einzelbild pruefen
exiftool photo.jpg | grep -E "GPS|Date|Time"
\end{lstlisting}

DMS zu Dezimalgrad umrechnen:
$$48^\circ\;4'\;35.00''\;\text{N} = 48 + \frac{4}{60} + \frac{35}{3600} = 48.0764^\circ$$

Fahrtroute rekonstruieren:
\begin{enumerate}
	\item GPS-Daten nach \texttt{DateTimeOriginal} sortieren
	\item Koordinaten in Dezimalgrad umrechnen
	\item Chronologisch in Karte einzeichnen
\end{enumerate}

Optional GPX-Track erzeugen:

\begin{lstlisting}[language=bash]
exiftool -p gpx.fmt /tmp/photos/*.jpg > /tmp/track.gpx
\end{lstlisting}

\section{YAFFS2 / Nexus Nandump}

Alte Android-Geräte (z.B. Nexus) nutzen YAFFS2.
Kein \texttt{mount} möglich, nur physische Analyse:

\begin{lstlisting}[language=bash]
# Dateisystem-Typ pruefen
fsstat nexus.nandump    # -> YAFFS

# Datenbanken suchen
fls -pr nexus.nandump | grep '\.db$' | grep -v '*'

# Geodaten extrahieren
icat nexus.nandump <INODE> > /tmp/CachedPosition.db
sqlitebrowser /tmp/CachedPosition.db
\end{lstlisting}
